\hypertarget{records}{We now introduce record types as an additional module of our framework. % by extending the grammar and the rules.
The basic intuition is that $\rectype{\Gamma}$ and $\recexp{\Gamma}$ construct record types and terms.}
We call a context \textbf{fully typed} resp. \textbf{defined} if all fields have a type resp. a definition.
In $\rectype{\Gamma}$, $\Gamma$ must be fully typed and may additionally contain defined fields.
In $\recexp{\Gamma}$, $\Gamma$ must be fully defined; the types are optional and usually omitted in practice.

Because we frequently need fully defined contexts, we introduce a notational convention for them: a context denoted by a \emph{lower} case letters like $\gamma$ is always fully defined.
In contrast, a context denoted by an \emph{upper} case letter like $\Gamma$ may have any number of types or definitions.

%\subsection{Records}

We extend our grammar as in \autoref{fig:recgrammar}.

\begin{figure}[ht]\centering%\vspace*{-1em}
  \begin{framed}
    \begin{tabular}{rcl@{\quad}l}	  
      \gray{$\Gamma$} & \gray{$::=$} & \gray{$\cdot \mid \Gamma, x[:T][:=T]$} &\\
      \gray{$T$} & \gray{$::=$} & \gray{$x \mid \type \mid \kind $} &\\
                     % & \gray{$\mid$} & \gray{$\prod_{x:T'}T \mid \lambda x:T'. T \mid T_1 T_2$} &\\
                      & $\mid$ &  $\rectype\Gamma \mid \recexp{\Gamma} \mid \recproj Tx$ & record types, terms, projections
    \end{tabular}
  \end{framed}%\vspace*{-.5em}
  \caption{Grammar for Records}\label{fig:recgrammar}\vspace*{-1em}
\end{figure}

\begin{remark}[Field Names and Substitution in Records]
Note that we use the same identifiers for variables in contexts and fields in records.
This allows reusing results about contexts when reasoning about and implementing records.
In particular, it immediately makes our records dependent, i.e., both in a record type and --- maybe surprisingly --- in a record term every variable $x$ may occur in subsequent fields.
In some sense, this makes $x$ bound in those fields.
However, record types are critically different from $\sigmasym$-types: we must be able to use $x$ in record projections, i.e., $x$ can \emph{not} be subject to $\alpha$-renaming.

As a consequence, capture-avoiding substitution is not always possible.
This is a well-known problem that is usually remedied by allowing every record to declare a name for itself (e.g., the keyword $\mathtt{this}$ in many object-oriented languages), which is used to disambiguate between record fields and a variable in the surrounding context (or fields in a surrounding record).
We gloss over this complication here and simply make substitution a partial function.
\end{remark}

Before stating the rules, we introduce a few critical auxiliary definition:

\begin{definition}[Substituting in a Record]\label{def:recsub1}\rm
We extend substitution $t \sub x{t'}$ to records:
    \begin{itemize}
    \item 
      $\rectype{x_1:T_1,\ldots,x_n:T_n}\sub yt\\\strut\qquad
        = \left\{\begin{array}{l}
                 \rectype{x_1:T_1\sub yt,\ldots,x_{i-1}:T_{i-1}\sub yt,x_i:T_i,\ldots, x_n:T_n} \text{if $y=x_i$}\\
                 \rectype{x_1:(T_1\sub yt),\ldots,x_n:(T_n \sub yt)}  \text{else}
               \end{array}
        \right.$ \\if none of the $x_i$ are free in $t$. Otherwise the substitution is undefined.
      \item 
        $\recexp{x_1:=t_1,\ldots,x_n:=t_n}\sub yt=
        \left\{\begin{array}{l}
                 \recexp{x_1:=t_1,\ldots,x_n:=t_n}\text{if $y\in\{x_1,\ldots,x_n\}$}\\
                 \recexp{x_1:=(t_1 \sub yt),\ldots,x_n:=(t_n \sub yt)} \text{else}
               \end{array}
       \right.$ \\if none of the $x_i$ are free in $t$. Otherwise the substitution is undefined.
      \item $(\recproj rx) \sub yt=\recproj{(r\sub yt)}x$.
      \end{itemize}
\end{definition}

\begin{definition}[Substituting with a Record]\label{def:recsub2}\rm
We write $t \sub r\Delta$ for the result of substituting any occurrence of a variable $x$ declared in $\Delta$ with $\recproj rx$

In the special case where $r=\recexp{\delta}$, we simply write $t[\delta]$ for $t\sub{\recexp{\delta}}\delta$, i.e., we have $t[x_1:=t_1,\ldots,x_n:=t_n]=t\sub{x_n}{t_n}\ldots\sub{x_1}{t_1}$.
\end{definition}

% % % % % % % % % % % % % % %

\begin{figure}[hbt]\centering
\begin{framed}\flushleft
\textbf{Formation:}
\[
	\nrule{\wfctx{\Gamma,\Delta} \quad \quad \Delta\text{ fully typed} \quad \umax\Delta\in\{\type,\kind\}}
		{\gjudg{\infertype{\rectype\Delta}{\umax\Delta} }}{}
\]
%where $\max \Delta$ is the maximal universe of all undefined fields in $\Delta$
%\[
%	\nrule{\gjudg{\judgeq R{\rectype{\Delta_1}}U} \quad \wfctx{\Gamma,\Delta_1+\Delta_2} \quad 
%		\max(\Delta_1+\Delta_2)\in\{\type,\kind\}}
%		{\gjudg{\infertype{\rectype{R \uplus \Delta_2}}{\max (\Delta_1+\Delta_2)}}}{\mathcal R_{F2}}
%\]

%\textbf{Type Elaboration:}
%\[
%	\nrule{\gjudg{\judgeq{R_1}{\rectype{[S_1\uplus]\Delta_1}}U_1} \quad \gjudg{\judgeq{R_2}{\rectype{[S_2\uplus]\Delta_2}}U_2} }
%		{\gjudg{\judgeq{R_1\uplus R_2}{\rectype{\mathtt{fl}(R_1)+\mathtt{fl}(R_2)}}{\max (\mathtt{fl}(R_1)+\mathtt{fl}(R_2))}}}
%		{\mathcal R_\uplus}
%\]

\textbf{Introduction:}
%\[
%	\nrule{\wfctx{\Gamma,\delta} \quad \delta\text{ fully defined}}
%		{\gjudg{\infertype{\recexp\delta}{\rectype\delta}}}{\mathcal R_{I1}}
%\] 
\[
	\irule{\wfctx{\Gamma,\Delta} \quad \delta\text{ fully defined}  \quad \Delta \text{ like $\delta$ but with all missing types inferred}}
 		{\gjudg{\infertype{\recexp \delta}{\rectype{\Delta}}}}%{\mathcal R_{I}}
 \]

\textbf{Elimination:}
\[
	\irule{\gjudg{\infertype r{\rectype{\Delta_1,x:T[:=t],\Delta_2}}}}
		{\gjudg{\infertype{\recproj rx}{T \sub r{\Delta_1}}}}%{\mathcal R_E}
\]

\textbf{Type checking:}
 \[
 	\irule{
 	\overbrace{\gjudg{\hastype{\recproj rx}{T\sub r\Delta}}}^\text{For $x:T[:=t]\in \Delta$ } \quad
 		\overbrace{\gjudg{\judgeq{\recproj rx}{t\sub r\Delta}{T\sub r\Delta}}}^\text{additionally for $x:T:=t\in \Delta$ }\quad
 	  \gjudg{\infertype rR}}
		{\gjudg{\hastype r\rectype{\Delta}}}%{\mathcal R_T}
\]

\textbf{Equality checking (extensionality):}
\[
	\irule{
		\overbrace{\gjudg{\judgeq{\recproj{r_1}x}{\recproj{r_2}x}{T\sub{r_1}\Delta}}}^\text{For every $(x:T)\in\Delta$}}
		{\gjudg{\judgeq{r_1}{r_2}{\rectype\Delta}}}%{\mathcal R_{Ext}}
\]

\textbf{Computation:}
\[
	\irule{\delta=(\delta_1,x[:T]:=t,\delta_2) \quad \gjudg{\infertype{\recexp{\delta}}{R}}}
		{\gjudg{\judgeqc{\recproj{\recexp{\delta}}x}{t[\delta_1]}}}%{\mathcal R_{C1}}
\quad
\irule{\gjudg{\infertype r\rectype{\Delta_1,x:T:=t,\Delta_2}}}
		{\gjudg{\judgeqc{\recproj rx}{t\sub r{\Delta_1}}} }%{\mathcal R_{C2}}
\]
\end{framed}%\vspace*{-.5em}
\caption{Rules for Records}\label{fig:rec:rules}%\vspace*{-1.em}
\end{figure}

Our rules for records are given in \autoref{fig:rec:rules}.
%Their roles are systematically similar to the rules for functions: three inference rules for the three constructors followed by a type and an equality checking rule for record types and the (in this case: two) computation rules.
We remark on a few subtleties below.
\begin{itemize}
\item The formation rule is partial in the sense that not every context defines a record type or kind.
This is because the universe of a record type must be as high as the universe of any undefined field to avoid inconsistencies.
% For example, $\max(a:\cn{nat})=\type$, $\max(a:\type)=\kind$ and $\max(a:\kind)$ is not defined.
If the presence of a hierarchy of universes (see \autoref{sec:lf:lfx:hierarchy}), we can turn every context into a record type.

\item The introduction rule infers the principal type of every record term.
Because we allow record types with defined fields, this is the singleton type containing only that record term.
This may seem awkward but does not present a problem in practice, where type checking is preferred over type inference anyway.
%In fact, it can make type inference more efficient because implementations may reuse the pointer to $\delta$ when synthesizing $\rectype{\Delta}$.

\item The elimination rule is straightforward, but it is worth noting that it is entirely parallel to the second computation rule.\footnote{Note that it does not matter how the fields of the record are split into $\Delta_1$ and $\Delta_2$ as long as $\Delta_1$ contains all fields that the declaration of $x$ depends on.}

\item The type checking rule has a surprising premise that $r$ must already be well-typed (against some type $R$).
Semantically, this assumption is necessary because we only check the presence of the fields required by $\rectype{\Delta}$ --- without the extra assumption, typing errors in any additional fields that $r$ might have could go undetected.
In practice, we implement the rule with an optimization: If $r$ is a variable or a function application, we can efficiently infer some type for it.
Otherwise, if $r=\recexp{\delta}$, some fields of $\delta$ have already been checked by the first premise, and we only need to check the remaining fields.
The order of premises matters in this case: we want to first use type checking for all fields for which $\rectype{\Delta}$ provides an expected type before resorting to type inference on the remaining fields.

\item In the equality checking rule, note that we only have to check equality at undefined fields --- the other fields are guaranteed to be equal by the assumption that $r_1$ and $r_2$ have type $\rectype{\Delta}$.

\item Like the type checking rule, the first computation rule needs	the premise that $r$ is well-typed to avoid reducing an ill-typed into a well-typed term.
In practice, our framework implements computation with a boolean flag that tracks whether the term to be simplified can be assumed to be well-typed or not; in the former case, this assumption can be skipped.

\item The second computation rule looks up the definition of a field in the type of the record.
Both computation rules can be seen uniformly as definition lookup rules --- in the first case the definition is given in the record, in the second case in its type.
\end{itemize}

\begin{remark}[Horizontal Subtyping and Equality]\label{remark:lf:lfx:records:horizontal}
As mentioned in \autoref{remark:prelim:mmt:horizontalsubtp}, \mmt's equality judgment could alternatively be formulated as an untyped equality $\judgequt t{t'}$. For our presentation of record types, however, the use of typed equality is critical.

For example, consider record values $r_1=\recexp{a:=1,b:=1}$ and $r_2=\recexp{a:=1,b:=2}$ as well as record types $R=\rectype{a:nat}$ and $S=\rectype{a:nat,b:nat}$. Due to \emph{horizontal subtyping}\footnote{Also called \emph{structural subtyping}; e.g., $\rectype{x:T,y:S}$ is a subtype of $\rectype{x:T}$. In contrast, e.g. the straightforward covariance rule for record types is called \emph{vertical} subtyping.}, we have $\subtp S R$ and thus both $\hastype{r_i}{S}$ and $\hastype{r_i}{R}$.
This has the advantage that the function $S\to R$ that throws away the field $b$ becomes the identity operation.
%For example, we can use the operation $(\lambda x:S.x):(S\to R)$ to convert from $S$ to $R$ instead of having to write a function $\lambda x:S.\recexp{a:=x.a}$ that throws away the field $b$.
Now our equality at record types behaves accordingly and checks only for the equality of those fields required by the type.
Thus, $\judgeq{r_1}{r_2}{R}$ is true whereas $\judgeq{r_1}{r_2}{S}$ is false, i.e., the equality of two terms may depend on the type at which they are compared.
While seemingly dangerous, this makes sense intuitively: $r_1$ can be replaced with $r_2$ in any context that expects an object of type $R$ because in such a context the field $b$, where $r_1$ and $r_2$ differ, is inaccessible.

Of course, this treatment of equality precludes downcasts: an operation that casts the equal terms $r_1:R$ and $r_2:R$ into the corresponding unequal terms of type $S$ would be inconsistent.
But such downcasts are still possible (and valuable) at the meta-level.
For example, a tactic $GroupSimp(G,x)$ that simplifies terms $x$ in a group $G$ can check if $G$ is commutative and in that case apply more simplification operations.

We will sometimes omit a type in the formal presentation of rules if convenient. The implementation in \mmt allows for equality rules to have an \emph{optional} argument for the type at which they operate, hence both typed an untyped equalities can be treated by the system.
\end{remark}

\begin{example}
  \autoref{fig:semilattices1} shows a record type of \defemph{Semilattices} (actually,
  this is a $\kind$ because it contains a $\type$ field) analogous to the grouping in \autoref{fig:rec:semi1} (using
  the usual encoding of axioms via judgments-as-types and higher-order abstract syntax for
  first-order logic).

%\begin{figure}[ht]%\centering
\begin{framed}
$\mathtt{Semilattice} := \leftrectype{\begin{array}{lll}
	U &&:\; \type \\
	\wedge &&:\; U \to U \to U \\
	\mathtt{assoc} &&:\; \ded \forall x,y,z:U.\; (x \wedge y) \wedge z \doteq x \wedge (y \wedge z) \\
	\mathtt{commutative} &&:\;\ldots \\
	\mathtt{idempotent} &&:\;\ldots \\
\end{array}}
\begin{array}{lll}
	U &&:\; \type \\
	\wedge &&:\; \lfarrow U{\lfarrow UU} \\
	\mathtt{assoc} &&:\; \ded \forall x,y,z:U.\; (x \wedge y) \wedge z \doteq x \wedge (y \wedge z) \\
	\mathtt{commutative} &&:\;\ldots \\
	\mathtt{idempotent} &&:\;\ldots \\
\end{array}
\rightrectype{\begin{array}{lll}
	U &&:\; \type \\
	\wedge &&:\; U \to U \to U \\
	\mathtt{assoc} &&:\; \ded \forall x,y,z:U.\; (x \wedge y) \wedge z \doteq x \wedge (y \wedge z) \\
	\mathtt{commutative} &&:\;\ldots \\
	\mathtt{idempotent} &&:\;\ldots \\
\end{array}}$
\end{framed}%\vspace*{-.5em}
\captionof{figure}{The (Record-)Kind of Semilattices }\label{fig:semilattices1}%\vspace*{-1em}
%\end{figure}

Then, given a record $r : \mathtt{Semilattice}$, we can form the record projection $\recproj r\wedge$, which has type $\lfarrow{\recproj rU}{\lfarrow{\recproj rU}{\recproj rU}}$ and $\recproj r{\mathtt{assoc}}$ yields a proof that $\recproj r\wedge$ is associative.
The intersection on sets forms a semilattice so (assuming we have proofs $\cap-\mathtt{assoc}$, $\cap\mathtt{-comm}$, $\cap\mathtt{-idem}$ with the corresponding types) we can give an instance of that type as
\[\csl : \mathtt{Semilattice} := \recexp{ U := \mathtt{Set}, \wedge := \cap, \mathtt{assoc} := \cap\mathtt{-assoc}, \ldots}\] 

%If we only care about the external perspective\ednote{stratified? I'm unclear with the terminology here}, we can define the usual order of a semilattice as:
%\[
%\mathtt{order} : \prod_{r : \mathtt{Semilattice}}r.U \to r.U \to \mathtt{bool} 
%	:= \lambda r: \mathtt{Semilattice}.\; \lambda x,y:r.U.\;x \doteq x\wedge y
%\]
%However, if we want to model more complex structures as record types, these definitions
%are not available \emph{within} the record type.\ednote{Use terminology from intro}

\end{example}

%As mentioned previously, record types naturally introduce a subtyping principle:
%\begin{definition}\rm\label{def:subtp}
%  \begin{enumerate}
%  \item Let $\gjudg{\hastype{T_1,T_2}U}$. Then $T_1$ is \textbf{a subtype of} $T_2$ in $\Gamma$ (in symbols: $\gjudg{\subtp{T_1}{T_2}}$) iff $\judg{\Gamma,x:T_1}{\hastype x{T_2}}$.
%  \item For $A=\{x_i :T_i \mid i\in I\}$ and $B=\{x_i :T_i \mid i\in J\}$ we let $\gjudg{A\dot\subseteq B}$ iff for every $x:T\in A$:
%    \begin{itemize}
%    \item $x:T\in B$ or
%    \item $x:T'\in B$ and $\gjudg{\subtp{T'}T}$
%    \end{itemize}
%  \end{enumerate}
%\end{definition}
%
%Furthermore, given the rules presented in this section, the following properties are easy to verify (proofs straight-forward and omitted for brevity):
%
%\begin{theorem}\label{thm:subtp}
%  \begin{enumerate}
%  \item \emph{Weakening, exchange, substitution} and \emph{subject reduction} hold for the dependently typed lambda calculus and are preserved when adding the rules for records.
%  \item All judgments are \emph{decidable}.
%  \item The following subtyping rule is admissible:	
%    \[ \nrule{\gjudg{\hastype{\rectype{x_i :T_i}_{i\in I},\rectype{x_i :T_i}_{i\in J}}U}\quad 
%	  	\gjudg{\{x_i :T_i \mid i\in I\} \dot\subseteq \{x_i :T_i \mid i\in J\}}}
%		{\gjudg{\subtp{\rectype{x_i :T_i}_{i\in J}}{\rectype{x_i :T_i}_{i\in I}}}}{\mathcal R'_{<:}} 
%    \]
%  \item \textbf{($\eta$-Rule)} For any well-formed record $r=\recexp{\ell_i:=t_i}_{1\leq i\leq n}:R$ we have 
%    \[\gjudg{\judgeq r{\recexp{x_i:=r.x_i}_{1\leq i\leq n}}R}.\]
%\end{enumerate}
%\end{theorem}

\begin{theorem}[Principal Types]\label{thm:lf:records:principal}
Our inference rules (in plain \LF) infer a principal type for each well-typed normal term.
\end{theorem}
\begin{proof}
	Let $\Gamma$ be a fixed well-typed context. We need to show that for any normal expression $t$ the inferred type is the most specific one, meaning if $\gjudg{\infertype tT}$, then for any $T'$ with $\gjudg{\hastype t{T'}}$ we have $\gjudg{\subtp T{T'}}$.
	
	If the only type checking rule applicable to a term $t$ is an inference rule, then the only way for $t$ to check against a type $T'$ which is not the inferred type $T$ is by first inferring $T$ and then checking $\gjudg{\subtp T{T'}}$, so in these cases the claim follows by default.
	
	By induction on the grammar:
	\begin{labeling}{$t=\lambda x:A.t'$}
		\item[$t=x$] and $x:T\in\Gamma$, then $\gjudg{\infertype tT}$, which is principal by default.
		\item[$t=\type$] then $\gjudg{\infertype t\kind}$, which is principal by default.
		\item[$t=\kind$] is untyped.
		\item[$t=\prod_{x:A}B$] then $\gjudg{\infertype tU}$, where $U\in\{\kind,\type\}$, both of which are principal by default.
		\item[$t=\lambda x:A.t'$] then for some $B:U$ we have $\judg{\Gamma,x:A}{\infertype {t'}B}$, which is principal by default.
		
		\item[$t=f a$] then $\gjudg{\infertype t{B[x/a]}}$, where $\gjudg{\infertype f{\prod_{x:A}B}}$ $\gjudg{\infertype a{A'}}$ are both principal types and for well-typedness $\gjudg{\subtp{A'}A}$ needs to hold. Since $t$ is normal, $f$ does not simplify to a lambda-expression and  $B[x/a]$ is principal by default.
		
		\item[$t=\rectype\Delta$] then $\gjudg{\infertype tU}$, where $U\in\{\kind,\type\}$, both of which are unique and hence principal.
		\item[$t=\recexp\delta$] then $\gjudg{\infertype r{\rectype \Delta}}$, wheren $\Delta$ contains the exact same variables, but with all types inferred (and by induction hypothesis principal). For $t$ to check against a type, it has to have the form $\rectype{\Delta'}$, hence we need to show that if $\gjudg{\hastype t{\rectype{\Delta'}}}$, then $\judg{\Gamma,r:\rectype\Delta}{\hastype r{\rectype{\Delta'}}}$.
		
		Consider the type checking rule for records in \autoref{fig:rec:rules} and let $x:T[:=d]\in\Delta'$. Since $\gjudg{\hastype t{\rectype{\Delta'}}}$, we have $\gjudg{\hastype{\recproj rx}T}$ (and if $x$ is defined in $\Delta'$ also $\gjudg{\judgeq{\recproj tx}dT}$) and since $\Delta$ is inferred from $t$ we have $x:T':=d$ in $\Delta$, where by hypothesis $T'$ is the principal type of $d$ and hence $\gjudg{\subtp{T'}T}$.
		
		As a result, $\judg{\Gamma,r:\rectype\Delta}{\infertype{\recproj rx}{T'}}$, therefore $\judg{\Gamma,r:\rectype\Delta}{\hastype{\recproj rx}T}$ and $\judg{\Gamma,r:\rectype\Delta}{\judgeq{\recproj rx}dT}$ and hence 
		$\judg{\Gamma,r:\rectype\Delta}{\hastype r{\rectype{\Delta'}}}$, which makes $\rectype\Delta$ the principal type of $t$.
		
		\item[$t=\recproj rx$] Since $t$ is normal, we have $\gjudg{\infertype{\recproj rx}T\sub r{\Delta_1}}$ for some $\gjudg{\infertype r{\rectype{\Delta_1,x:T,\Delta_2}}}$. Since the latter is by hypothesis the principal type of $r$, for $t$ to typecheck against some $T'$ it needs to be the case that $r$ type checks against some $R=\rectype{\Delta_1',x:T',\Delta_2'}$ and $\gjudg{\subtp T{T'}}$ holds by the principal type of $r$.
	\end{labeling}
\end{proof}

In analogy to function types, we can derive the subtyping properties of record types.
We introduce context subsumption and then combine horizontal and vertical subtyping in a single statement.

\begin{definition}[Context Subsumption]\label{def:subsume}\rm
\hypertarget{subsume}{For two fully typed contexts $\Delta_i$ we write $\gjudg{\subsume{\Delta_1}{\Delta_2}}$ iff for every declaration $x:T[:=t]$ in $\Delta_1$ there is a declaration $x:T'[:=t']$ in $\Delta_2$ such that}
  \begin{compactitem}
  \item $\gjudg{\subtp{T'}T}$ and
  \item if $t$ is present, then so is $t'$ and $\gjudg{\judgeq{t}{t'}{T}}$
  \end{compactitem}
\end{definition}
\ednote{annoyingly, the subtyping/equality judgments need a different context: $\Delta_1^x\oplus\Delta_2^x$ where $^x$ means "the prefix up until and excluding $x$".}
%It seems we need to use $\oplus$ already and make this definition mutually recursive with a proof that subsumption implies that $\subsume{\Delta_1}{\Delta_2}$ implies that $\Delta_1^x\oplus\Delta_2^x$ is defined\\DM: I think the ``excluding $x$'' part is obvious to any reader, and the ``up until''-part is relatively benign, since fields in a record[type] can't depend on latter fields.}

Intuitively, $\subsume{\Delta_1}{\Delta_2}$ means that everything of $\Delta_1$ is also in $\Delta_2$.
That yields:

\begin{theorem}[Record Subtyping]\label{thm:recsub}
The following rule is derivable:

	\[ \irule{
	    %\gjudg{\hastype{\rectype{\Delta_1}}U}\quad 
    	%\gjudg{\hastype{\rectype{\Delta_2}}U}\quad 
			\gjudg{\subsume{\Delta_1}{\Delta_2}}}
			{\gjudg{\subtp{\rectype{\Delta_2}}{\rectype{\Delta_1}}}}%{\mathcal R_{<:}} 
		\]
\end{theorem}
\begin{proof}
	Assume $\gjudg{\subsume{\Delta_1}{\Delta_2}}$. We need to show $\judg{\Gamma,r:\rectype{\Delta_2}}{\hastype r{\rectype{\Delta_1}}}$. By the type checking rule in \autoref{fig:rec:rules}, for any $x:T[:=t]\in\Delta_1$, we need to show that $\judg{\Gamma,r:\rectype{\Delta_2}}{\hastype{\recproj rx}T}$ (and if applicable $\judg{\Gamma,r:\rectype{\Delta_2}}{\judgeq{\recproj rx}tT}$).
	
	By definition of $\subsume{\Delta_1}{\Delta_2}$, since $x:T[:=t]\in\Delta_1$, we have $x:T'[:=t]\in\Delta_2$ and $\gjudg{\subtp{T'}T}$, and if $x$ is defined in $\Delta_1$ the required equality holds as well, so the type checking rule proves $\judg{\Gamma,r:\rectype{\Delta_2}}{\hastype r{\rectype{\Delta_1}}}$ and the claim follows.
\end{proof}

%
%
%
%The intuition being, that if a record type $\rectype{\Delta}$ has a defined field $\ell_k:T_k:=t_k$, an instance $r$ of that type need not have that field defined again in the record itself. Hence $\mathcal R'_T$ only checks for the undefined fields in $\rectype\Delta$, but by $\mathcal R_E$ the projection $r.\ell_k$ is still well-formed and by $\mathcal R'_{C2}$ simplifies to the definition $t_k$ given in $\rectype\Delta$. 
%
%Since all judgments regarding records reduce to judgments in basic LF, adding those rules preserves the validity of weakening, exchange, substitution and subject reduction. Notably, $\mathcal R'_{C2}$ breaks decidability, since checking the conclusion requires finding \emph{some} record type $\rectype{\Delta_{T[D]}}$ with a defined field $\ell_k:T_k$ (in an implementation, the premise would be restricted to explicitly declared record types in the current context, of course).
%\ednote{mention problems if multiple records with the same field/type but different definitions exist}
%

%
%However this leads to problems if we assume the usual definition for subtyping and want extensionality of types (i.e. $\Gamma\vdash T_1<:T_2$ and $\Gamma\vdash T_2<:T_1$ implies $\Gamma\vdash T_1=T_2 : U$) and transitivity of subtyping to hold:

%
%By definition, reflexivity, transitivity and antisymmetry for subtyping follow immediately. Now the following subtyping rule becomes admissible:
%
%\begin{thm}
%	The following subtyping rule is admissible:	
%	\[ \dfrac{\Gamma\vdash \rectype{\Delta_1},\rectype{\Delta_2}:U\quad \Delta_2\subseteq \Delta_1}{\Gamma\vdash \rectype{\Delta_1} <: \rectype{\Delta_2}}\quad\mathcal R_{<:}
%	\]
%\end{thm}
%\begin{proof}
%	Follows immediately from $\mathcal R'_T$, since if $\Gamma\vdash r:\Delta_1$, then all projections $r.\ell$ for $\ell\in\Delta_2$ are well-formed and have the right types to infer $\Gamma\vdash r:\Delta_2$.
%\end{proof}
%
%\begin{thm}
%	Given the rules in Figure \ref{fig:naiverules}, all records types are equal.
%\end{thm}
%\begin{proof}
%	We'll show that every record type $\rectype\Delta$ is equal to the empty record type $\rectype\emptyset$ inductively over the length of $\Delta$. The case $\Delta=\emptyset$ is trivial, so let $\Delta=\Delta';x:T$ (wlog $x$ undefined in $\Delta$) and assume $\rectype{\Delta'}=\rectype\emptyset$. 
%	
%	By $\mathcal R_{<:}$ we have $\rectype\Delta<:\rectype\emptyset=\rectype{\Delta'}$. Conversely, let $r:\rectype{\Delta'}$, $t:T$ and consider $\rectype{\Delta''}=\rectype{\Delta';x:T:=t}$. By $\mathcal R'_T$ we have $r:\rectype{\Delta''}$ and hence $\rectype{\Delta'}<:\rectype{\Delta''}$. But since for any $r:\rectype{\Delta''}$ the projection $r.x$ is (by $\mathcal R'_{C2}$) well-formed, equal to $t$ and hence has type $T$, we also get $r:\rectype\Delta$, hence $\rectype{\Delta''}<:\rectype{\Delta}$ and by transitivity $\rectype{\Delta'}<:\rectype\Delta$.\qed
%\end{proof}
%
%Technically, the above proof is not as devastating on an implementation level, since for $r.x$ to be a well-formed expression, the projection has to be explicitly checked in the context of $r:\rectype{\Delta''}$, which is not the inferable type from $\mathcal R_I$. However, it's still unsatisfying on a theoretical level and prevents the explicit implementation of admissible subtyping rules that are desirable for efficiency.
%
%To fix the above problem, we'll adapt $\mathcal R'_I, \mathcal R'_T, \mathcal R'_{C1}$ and $\mathcal R'_{C2}$ and allow for records and record types to have a designated field $\self$, as in Figure \ref{fig:recrules}: 
%\begin{itemize}
%\item In a record type $T$, a field $\self : R$ is purely for convenience reasons and signifies that $T$ extends $R$ by the explicitly declared fields in $T$ (in particular $T<:R$ is enforced). Consequently, we consider $\rectype{\self : \rectype{\Delta_1};\Delta_2}$ to be a mere abbreviation for $\rectype{\Delta_1;\Delta_2}$, where $\Delta_1,\Delta_2$ need to be compatible in the sense that if $\Delta'$ is the subcontext of all variables in $\Delta_2$ occurring in $\Delta_1$, then $\Delta'\subseteq\Delta_1$ needs to hold.
%\item In a record $r$, a field $\self : R$ serves as a type annotation by ``importing'' all defined fields in the record type $R$ and introducing the obligation to implement all non-defined fields. Consequently, if the record $r$ in the above proof does not explicitly have the field $\self : \rectype{\Delta''}$, it does not typecheck against $\rectype{\Delta''}$, hence the subtyping judgment $\rectype{\Delta'}<:\rectype{\Delta''}$ does not hold.
%\end{itemize}
%
%\begin{figure}]
%\fbox{\begin{minipage}{\textwidth}
%
%\textbf{Introduction (Type Inference):}
%\\ Let $\Delta'$ be $\Delta$ restricted to the variables not defined in $\Delta_{T[D]}$.
%
%\[\dfrac{\Gamma;\Delta\vdash\cdot \quad \self\notin\Delta}
%{\Gamma\vdash\recexp{\Delta_D} \overset{\rightsquigarrow}{:} \rectype{\Delta_T}}\quad\mathcal R_{I1} 
%\qquad
%\dfrac{\Gamma;\self : \rectype{\Delta_{T[D]}} ;\Delta'\vdash\cdot}
%{\Gamma\vdash\recexp{\self : \Delta_{T[D]};\Delta'} \overset{\rightsquigarrow}{:} \rectype{\Delta_{T[D]} } }\quad\mathcal R_{I2}
%\]
%
%\textbf{Typing:}
%\[
%\dfrac{\Gamma\vdash\rectype{\Delta_{T[D]}}:U \quad\overbrace{\Gamma\vdash r.\ell_i:T_i[r]}^\text{For $\ell_i\in Dec\left(\Delta_{T[D]}\right)$ } \quad \overbrace{\Gamma\vdash r.\ell_i=t_i[r]:T_i}^\text{For $\ell_i\notin Dec\left(\Delta_{T[D]}\right)$} }
%{\Gamma\vdash r : \rectype{\Delta_{T[D]}}} \quad \mathcal R_T
%\]
%
%\textbf{Computation:} \\ For $r:=\recexp{\Delta_D}$:
%\[\dfrac{\Gamma\vdash r:\rectype{\Delta_T} \quad \ell_k\in \Delta_D }
%{\Gamma\vdash r.\ell_k=t_k\,[r]:T_k\,[r]}\quad\mathcal R_{C1}\]
%For $r:=\recexp{\self : \Delta_{T[D]}; R'}$:
%\[\dfrac{\Gamma\vdash\rectype R:U \quad \Gamma\vdash r:\rectype R \quad \ell_k\text{ defined in }\Delta_{T[D]} \quad \ell_k\notin R' }
%{\Gamma\vdash r.\ell_k=t_k\,[r]:T_k\,[r]}\quad\mathcal R_{C2}\]
%
%\end{minipage}}
%\caption{The updated rules for records}
%\label{fig:recrules}
%\end{figure}
%
% Given the new rules, decidability is preserved since all premises either require arbitrary record types ($\mathcal R_{C2}$, where the first two premises are only necessary to guarantee well-formedness of $r$) where we can always fall back to the inferable ones, or types uniquely inferable from the conclusions ($\mathcal R_{C1}$, $\mathcal R_T$).
%
%Furthermore, $\mathcal R_{<:}$ remains admissible. Given our definition of $\subseteq$ on contexts, we can give more specific rules (for efficiency reasons):
%
%\begin{thm}
%	The following rules are admissible:
%	\[
%	\dfrac{\Gamma;\Delta_1\vdash t:A\quad\Gamma;\Delta_1;x:A:=t;\Delta_2\vdash\cdot}
%	{\Gamma\vdash\rectype{\Delta_1;x:A:=t;\Delta_2} <: \rectype{\Delta_1;x:A;\Delta_2}} \quad \mathcal R_{<:D}
%	\]
%	
%	\[
%	\dfrac{\Gamma;\Delta\vdash \cdot \quad \Gamma\vdash A:U}
%	{\Gamma\vdash\rectype{\Delta;x:A} <: \rectype{\Delta}} \quad \mathcal R_{<:+}
%	\]
%	
%	
%	\[
%	\dfrac{\Gamma;\Delta_1\vdash A <: A'\quad\Gamma;\Delta_1;x:A;\Delta_2\vdash\cdot}
%	{\Gamma\vdash\rectype{\Delta_1;x:A;\Delta_2} <: \rectype{\Delta_1;x:A';\Delta_2}} \quad \mathcal R_{<:<:}
%	\]
%\end{thm}
%
%In the presence of subtyping, uniqueness of types is naturally lost. However, we can show the existence of \emph{principal types}, which are sufficient in most implementation-related situations:
%
%\begin{thm}
%	Every well-typed term $t$ has a \defemph{principal type} $T$ in the sense that $\Gamma\vdash t:T$ and whenever $\Gamma\vdash t:T'$, then also $\Gamma\vdash T<:T'$.
%
%	Furthermore, the principal type of a record term is exactly the type inferred by rule $\mathcal R_I$.
%\end{thm}
%\begin{proof} \ednote{update proof}
%	Keep in mind that subtyping only ever occurs due to records. We proceed by induction on $t$: 
%	
%	\begin{itemize}	
%		\item The cases $t\in\{s,\type,\kind,v\}$ are trivial.
%		\item If $t=\prod_{v:t_1}t_2$, then the type of $t$ is either $\type$ or $\kind$ and hence unique.
%		\item If $t=\lambda v:t_1.t_2$ and $\Gamma\vdash t:T$, then by induction hypothesis $t_2$ has a principal type $T_2$ and $T$ has the form $\prod_{v:t'_1}T_2'$. We claim the principal type of $t$ is $\prod_{v:t_1}T_2$: 
%				
%			By variance we need to show $\Gamma\vdash t_1'<:t_1$ and $\Gamma\vdash T_2 <: T_2'$. The former follows immediately from $\Gamma\vdash t:T$. For the latter, since $t:T$, it needs to be the case that $\Gamma,v:t_1\vdash t_2:T_2'$, hence the claim follows by principality of $T_2$. 
%				
%		\item If $t=t_1 t_2$, then $t_1$ has a principal type of the form $\prod_{v:A}B$ and hence $t$ has the unique type $B$.
%		\item If $t=\recexp{\ell_i:=t_i}_{i\in I}$ is a record, then clearly $\Gamma\vdash t:T_t=:\rectype{\ell_i:T_i}_{i\in I}$. For any other record type $T=\rectype{\ell_i:T'_i}_{i\in I'}$ with $\Gamma\vdash t:T$, by $\mathcal R_T$ (and $\mathcal R_C$) we have $I'\subseteq I$, and since eliminating fields in record types leads to (either incomparable types or) supertypes, we can wlog assume $I=I'$. By $\mathcal R_T$ we have $\Gamma\vdash t_i:T_i$ and $\Gamma\vdash t_i:T'_i$ for all $i\in I$. By induction hypothesis the $T_i$ are all principle, hence $\Gamma\vdash T_i<:T'_i$ and hence $\Gamma\vdash T_t <: T$.
%	\end{itemize}
%\end{proof}

\paragraph{Merging Records}\label{sec:merging}

We introduce an advanced operation on records, which proves critical for both convenience and performance:
Theories can easily become very large containing hundreds or even thousands of declarations.
If we want to treat theories as record types, we need to be able to build big records from smaller ones without exploding them into long lists.
Therefore, we introduce an explicit merge operator $\recmergesym$ on both record types and terms.

\hypertarget{merge}{In the grammar, this is a single production for terms:}
\[T ::= \recmerge TT\]
The intended meaning of $\recmergesym$ is given by the following definition:
\begin{definition}[Merging Contexts]\label{def:recmerge}\rm
Given a context $\Delta$ and a (not necessarily well-typed) context $E$, we define a partial function $\contmerge\Delta E$ as follows:
    \begin{compactitem}
    \item $\contmerge\cdot E \;=\; E$
    \item If $\Delta=d,\Delta_0$ where $d$ is a single declaration for a variable $x$:
     \begin{compactitem}
      \item if $x$ is not declared in $E$:
        $\contmerge{(d,\Delta_0)}E \;=\; d,\;(\contmerge{\Delta_0} E)$
      \item if $E\;=\;E_0,e,E_1$ where $e$ is a single declaration for a variable $x$:
       \begin{compactitem}
        \item if a variable in $E_0$ is also declared in $\Delta_0$: $\contmerge\Delta E$ is undefined,
        \item if $d$ and $e$ have unequal types or unequal definitions: $\contmerge\Delta E$ is undefined\footnotemark,
        \item otherwise, $\contmerge{(d,\Delta_0)}{(E_0,e,E_1)}=E_0,m,(\Delta_0,E_1)$ where $m$ arises by merging $d$ and $e$.
       \end{compactitem}
     \end{compactitem}
    \end{compactitem}
\end{definition}
\footnotetext{It is possible and important in practice to also define $\contmerge\Delta E$ when the types/definitions in $d$ and $e$ are provably equal.
We omit that here for simplicity.}
Note that $\contmergesym$ is an asymmetric operator: While $\Delta$ must be well-typed (relative to some ambient context), $E$ may refer to the names of $\Delta$ and is therefore not necessarily well-typed on its own.

%To fortify our intuition, we prove:
%\begin{theorem}\label{thm:merge}
%Given fixed $\Gamma$, if $\Delta\oplus E$ is defined, then $\wfctx{\Gamma,\Delta\oplus E}$ and $\gjudg{\subsume{\Delta}{\Delta\oplus E}}$.
%%		\[
%%			\irule{\gjudg{\subtp R{R_1}} \quad \gjudg{\subtp R{R_2}} \quad \gjudg{\hastype{R_1+R_2}U}}{\gjudg{\subtp R{R_1+R_2}}}
%%		\]
%\end{theorem}
%\begin{proof}
%	See appendix.
%\end{proof}

We do not define the semantics of $\recmergesym$ via inference and checking rules.
Instead, we give equality rules that directly expand $\recmergesym$ into $\contmergesym$ when possible:

\[
	\nrule{\wfctx{\Gamma,(\contmerge{\Delta_1}{\Delta_2})}}
		{\gjudg{\judgeqc{\recmerge{\rectype{\Delta_1}}{\rectype{\Delta_2}}}{\rectype{\contmerge{\Delta_1}{\Delta_2}}}{}}}
		{}
\quad
	\nrule{\wfctx{\Gamma,(\contmerge{\delta_1}{\delta_2})}}
		{\gjudg{\judgeqc{\recmerge{\recexp{\delta_1}}{\recexp{\delta_2}}}{\recexp{\contmerge{\delta_1}{\delta_2}}}{}}}
		{}
\]\[ %\quad
	\nrule{\wfctx{\Gamma,(\contmerge\Delta\delta)}}
		{\gjudg{\judgeqc{\recmerge{\rectype{\Delta}}{\recexp{\delta}}}{\recexp{\contmerge\Delta\delta}}{}}}
		{}
\]
In implementations some straightforward optimizations are needed to verify the premises of these rules efficiently; we omit that here for simplicity.
For example, merges of well-typed records with disjoint field names are always well-typed, but e.g., $\recmerge{\rectype{x:nat}}{\rectype{x:bool}}$ is not well-typed even though both arguments are.

In practice, we want to avoid using the computation rules for $\recmergesym$ whenever possible.
Therefore, we prove admissible rules (i.e., rules that can be added without changing the set of derivable judgments) that we use preferentially:
\begin{theorem}\label{thm:plusrules}
\begin{itemize}
	\item If $R_1$, $R_2$, and $\recmerge{R_1}{R_2}$ are well-typed record types, then $\recmerge{R_1}{R_2}$ is the greatest lower bound with respect to subtyping of $R_1$ and $R_2$.
In particular, $\gjudg{\hastype r{\recmerge{R_1}{R_2}}}$ iff $\gjudg{\hastype r{R_1}}$ and $\gjudg{\hastype r{R_2}}$

	\item If $\gjudg{\hastype{r_i}{R_i}}$ and $\recmerge{r_1}{r_2}$ is well-typed, then $\gjudg{\hastype{\recmerge{r_1}{r_2}}{\recmerge{R_1}{R_2}}}$
	\end{itemize}
\end{theorem}
\begin{proof}
	\begin{itemize}
		\item If $R_1=\rectype{\Delta_1}$ and $R_2=\rectype{\Delta_2}$ are well-typed record types, then
			$\judgeqc{\recmerge{R_1}{R_2}}{\rectype{\contmerge{\Delta_1}{\Delta_2}}}$. By definition of $\contmergesym$, any $r:\recmerge{R_1}{R_2}$ hence has all fields defined that are required by both $\rectype{\Delta_1}$ and $\rectype{\Delta_2}$ and the other way around.
		\item By the rules for $\recmergesym$, for $r_1=\recexp{\delta_1}$ and $\recexp{\delta_2}$ we get $\recmerge{r_1}{r_2}=\recexp{\contmerge{\delta_1}{\delta_2}}$ and if $R_1=\rectype{\Delta_1}$ and $R_2=\rectype{\Delta_2}$, then (since $\gjudg{\hastype{r_i}{R_i}}$)  we have $\subsume{\Delta_1}{\delta_1}$ and $\subsume{\Delta_2}{\delta_2}$.
		
		As can easily be verified, it follows that $\subsume{\contmerge{\Delta_1}{\Delta_2}}{\contmerge{\delta_1}{\delta_2}}$, and hence $\infertype{\recexp{\contmerge{\delta_1}{\delta_2}}}{\rectype{\contmerge{\delta_1}{\delta_2}}}<:\rectype{\contmerge{\Delta_1}{\Delta_2}}=\recmerge{R_1}{R_2}$.
	\end{itemize}
\end{proof}

Inspecting the type checking rule in \autoref{fig:rec:rules}, we see that a record $r$ of type $\rectype\Delta$ must repeat all defined fields of $\Delta$.
This makes sense conceptually but would be a major inconvenience in practice.
The merging operator solves this problem elegantly as we see in the following example:

\begin{example} 
  Continuing our running example, we can now define a type of semilattices \emph{with order} (and all associated axioms) as in \autoref{fig:semi2}.
%\begin{figure}[ht]\centering
\begin{framed}
$\mathtt{SemilatticeOrder} := \mathtt{Semilattice} \recmergesym \leftrectype{\begin{array}{lll}
%	\gray U &&\gray:\; \gray\type \\
%	\gray \wedge &&\gray:\; \gray{U \to U \to U} \\
%	\gray{\mathtt{assoc}} &&\gray:\; \gray{\vdash \forall x,y,z:U.\; (x \wedge y) \wedge z \doteq x \wedge (y \wedge z)} \\
%	\gray{\mathtt{commutative}} &&\gray:\;\gray\ldots \\
%	\gray{\mathtt{idempotent}} &&\gray:\;\gray\ldots \\
          \leq &&:\; \lfarrow U{ \lfarrow UU} := \lflambda{x,y:U}{x \doteq x\wedge y}\\
           \mathtt{refl} &&:\; \vdash \forall a:U.\; a \leq a := \text{(proof)}\\
           \ldots \\
 \end{array}} 
 \begin{array}{lll}
%	\gray U &&\gray:\; \gray\type \\
%	\gray \wedge &&\gray:\; \gray{U \to U \to U} \\
%	\gray{\mathtt{assoc}} &&\gray:\; \gray{\vdash \forall x,y,z:U.\; (x \wedge y) \wedge z \doteq x \wedge (y \wedge z)} \\
%	\gray{\mathtt{commutative}} &&\gray:\;\gray\ldots \\
%	\gray{\mathtt{idempotent}} &&\gray:\;\gray\ldots \\
          \leq &&:\; \lfarrow U{ \lfarrow UU} := \lflambda{x,y:U}{x \doteq x\wedge y}\\
           \mathtt{refl} &&:\; \ded \forall a:U.\; a \leq a := \text{(proof)}\\
           \ldots \\
 \end{array}
                        \rightrectype{\begin{array}{lll}
%	\gray U &&\gray:\; \gray\type \\
%	\gray \wedge &&\gray:\; \gray{U \to U \to U} \\
%	\gray{\mathtt{assoc}} &&\gray:\; \gray{\vdash \forall x,y,z:U.\; (x \wedge y) \wedge z \doteq x \wedge (y \wedge z)} \\
%	\gray{\mathtt{commutative}} &&\gray:\;\gray\ldots \\
%	\gray{\mathtt{idempotent}} &&\gray:\;\gray\ldots \\
          \leq &&:\; \lfarrow U{ \lfarrow UU} := \lflambda{x,y:U}{x \doteq x\wedge y}\\
           \mathtt{refl} &&:\; \vdash \forall a:U.\; a \leq a := \text{(proof)}\\
           \ldots \\
 \end{array}}$\\
$\cslo  :=  \recmerge{\mathtt{SemilatticeOrder}}\csl$
\end{framed}%\vspace*{-.5em}
\captionof{figure}{Running Example}\label{fig:semi2}%\vspace*{-1em}
%\end{figure}

Now the explicit merging in the type $\mathtt{SemilatticeOrder}$ allows the projection $\recproj\cslo\leq$, which is equal to $\lambda x,y:(\recproj\cslo U)\;.\;(x \doteq x (\recproj\cslo\wedge) y)$ and $\recproj\cslo{\mathtt{refl}}$ yields a proof that this order is reflexive -- without needing to define the order or prove the axiom anew for the specific instance $\csl$.
\end{example}
